\documentclass[10pt,a4paper]{amsart}
\usepackage[margin=19mm]{geometry}
\usepackage{amsmath,amssymb,mathtools}
\usepackage[scr=rsfs,cal=euler]{mathalfa}
\usepackage{xfrac}
\usepackage[unicode,hidelinks,bookmarks=false]{hyperref}
\newcommand{\ie}{i.e.\ }
\newcommand{\cf}{cf.\ }
\newcommand{\resp}{respectively\ }
\newcommand{\Subsection}{\subsection}
\providecommand{\nobreakdash}{\nobreak\hbox{-}}
\setlength{\emergencystretch}{2em}
\multlinegap0pt
\usepackage{amssymb}
\newtheorem{theorem}{Theorem}
\newtheorem{lemma}{Lemma}
\newcommand{\Bn}{\mathbb B^n}
\newcommand{\Dh}{\Delta_h}
\title{A Converse Mean-Value Property}
\author{Mat\v{e}j Morav\'ik, El Hassan Youssfi}
\date{Source: arXiv:2606.22199v1 (CC BY 4.0)}
\begin{document}
\maketitle

\noindent\emph{Source and licence.} A Converse Mean-Value Property. Version arXiv:2606.22199v1 (CC BY 4.0), distributed by arXiv under the Creative Commons Attribution 4.0 International licence, http://creativecommons.org/licenses/by/4.0/, as recorded in the arXiv licence metadata for this exact version and shown on the abstract page https://arxiv.org/abs/2606.22199v1. Full licence notice: \texttt{licenses/arxiv-2606.22199-CC-BY-4.0.txt}.\par\smallskip

\noindent\textit{Editorial scope.} Counted unit: the article's lemma Regularity (source0.tex lines 470--477). The article's complete original proof of it is delivered with it (source0.tex lines 479--553, one \texttt{proof} environment of 75 lines). No mathematics is written, repaired or re-derived: the statement and the proof are the article's own lines, in the article's own notation, and no retained line is substituted or re-flowed.  Retained uncounted as the source-local context the proof needs: lines 209--244; lines 248--303; lines 317--334.  Nothing was omitted from any retained span, so the package carries no omission record.  No retained line contains a citation, so the wrapper carries no bibliography and no bibliography entry.  No retained line contains a figure environment, an image-inclusion command or drawing code, so no image is delivered and the package declares no asset dependency.  Editorial formatting only, in addition: an AMS \texttt{amsart} preamble that restates the article's own declarations and macros used by the retained lines, namely the environments \texttt{theorem}, \texttt{lemma} on separate counters, together with the standard packages those lines need.  The retained environments print in this document as Theorem 1, Lemma 1 in order of appearance, whereas the article prints the same items with the numbering of its own full document.  The complete native source of the article as distributed in the arXiv e-print for this exact version is bundled unchanged as \texttt{sources/203331/source0.tex}.
 Our main result is the follwing:
\begin{theorem}\label{main}
Suppose that \(n\ge2\) and let \(\Omega\subset \Bn\) be an open connected set  containing the origin $0$ such that
\begin{eqnarray}\label{finiteVol}
\nu(\Omega)<+\infty.
\end{eqnarray}
Assume that for every compact set \(K\Subset\mathbb B^n\),
\begin{eqnarray}\label{C1cond}
\sup_{z\in K}
\int_\Omega |\xi-z|^{1-n}\,d\nu(\xi)
<+\infty,
\end{eqnarray}
\begin{eqnarray}\label{C1bcond}
\lim_{r\to1^-} (1-r^2)^{n-2}
\sup_{z\in\Omega\cap\{|z|=r\}}
\int_\Omega |\xi-z|^{1-n}\,d\nu(\xi)=0
\end{eqnarray}
and
\begin{eqnarray}\label{Topcond}
\partial\Omega\cap\mathbb B^n =\partial\overline{\Omega}\cap\mathbb B^n.
\end{eqnarray}
%%%%%%%%
%%%%%%
%%%%%%%%%%%%
%%%%%%%%%
Suppose that, for every $\nu$-integrable  \(H\)-harmonic function \(f\) in a neighbourhood of
\(\overline{\Omega} \cap\mathbb B^n\), one has
\[
\frac1{\nu(\Omega)}\int_\Omega f(x)\,d\nu(x)=f(0).
\]
Then \(\Omega\) is a hyperbolic ball centered at \(0\). Equivalently, there
exists \(0<R<1\) such that
\[
\Omega=\{x\in\Bn:\ |x|<R\}.
\]
\end{theorem}

Let \(G(\xi,z)\) be the Green function associated with \(\Dh\). We use the
sign convention
\[
\Delta_{h,z} G(\xi,z)=\delta_\xi(z)
\]
in the sense of distributions. With this convention, the Green potential
\[
([G]\Psi)(z):=\int_{\Bn}\Psi(\xi)G(\xi,z)\,d\nu(\xi)
\]
satisfies
\[
\Dh ([ G]\Psi)=\Psi,
\]
for all $C^\infty$-compactly supported functions $\Psi$ in $\Bn.$
The Green function is symmetric:
\[
G(\xi,z)=G(z,\xi),
\]
and, for fixed \(z \in \mathbb B^n\), the function
\[
\xi\longmapsto G(\xi,z)
\]
is \(H\)-harmonic away from \(\xi=z\).
In particular, if \(z\notin\overline{\Omega}\), then
\[
\xi\longmapsto G(\xi,z)
\]
is a $\nu$-integrable \(H\)-harmonic in a neighbourhood of \(\overline{\Omega}\).
Put
\[
c=\nu(\Omega).
\]
The mean-value hypothesis says
\[
\int_\Omega f(\xi)\,d\nu(\xi)=c f(0)
\]
for every $\nu$-integrable \(H\)-harmonic function \(f\) in a neighborhoord of  \(\overline{\Omega}\cap\mathbb B^n\).
Taking
\[
f(\xi)=G(\xi,z),
\qquad z\notin\overline{\Omega},
\]
we obtain
\[
\int_\Omega G(\xi,z)\,d\nu(\xi)=cG(0,z),
\qquad z\notin\overline{\Omega}.
\]
Define
\begin{eqnarray}\label{h}
h(z):=\int_\Omega G(\xi,z)\,d\nu(\xi).
\end{eqnarray}
Then
\begin{eqnarray}\label{hG}
h(z)=cG(0,z),
\qquad z\notin\overline{\Omega}.
\end{eqnarray}

\begin{lemma}\label{gradient}
 Then there exists a constant \(C>0\) such that
\[
|\nabla_zG(\xi,z)|
\le
C(1-|z|^2)^{n-2}|\xi-z|^{1-n},
\qquad
\xi,z\in\Bn,\quad \xi\ne z.
\]
Consequently,
\[
|T_{ij,z}G(\xi,z)|
\le
C(1-|z|^2)^{n-2}|\xi-z|^{1-n},
\qquad
\xi,z\in\Bn,\quad \xi\ne z.
\]
\end{lemma}

\begin{lemma}\label{Regularity} Under (\ref{finiteVol}) and (\ref{C1cond}), the function $h$ given by (\ref{h})
satisfies
\begin{enumerate}
\item \(h\in C^1(\Bn)\);
\item \(\Delta_h h=\chi_\Omega\) in \(\mathcal D'(\Bn)\);
\item \(\Delta_h h=1\) classically in \(\Omega\).
\end{enumerate}
\end{lemma}

\begin{proof} 
Fix a compact set \(K\Subset\Bn\). Lemma \ref{gradient} implies that
\[
|\nabla_zG(\xi,z)|
\le C_K |\xi-z|^{1-n},
\qquad z\in K.
\]
By assumption,
\[
\sup_{z\in K}
\int_\Omega |\xi-z|^{1-n}\,d\nu(\xi)<+\infty.
\]
Hence
\[
\sup_{z\in K}
\int_\Omega |\nabla_zG(\xi,z)|\,d\nu(\xi)<+\infty.
\]
Therefore differentiation under the integral sign is justified and
\[
\partial_{z_j}h(z)
=
\int_\Omega \partial_{z_j}G(\xi,z)\,d\nu(\xi).
\]

The same domination allows the use of dominated convergence, proving that the
first derivatives are continuous. Thus
\[
h\in C^1(\Bn).
\]
Let \(\varphi\in C_c^\infty(\mathbb B^n)\). Using Fubini's theorem,
\[
\langle \Delta_h h,\varphi\rangle
=
\langle h,\Delta_h\varphi\rangle
=
\int_\Omega
\left(
\int_{\Bn}
G(\xi,z)\Delta_h\varphi(z)\,d\nu(z)
\right)
d\nu(\xi).
\]
Since
\[
\Delta_{h,z}G(\xi,z)=\delta_\xi(z),
\]
we obtain
\[
\int_{\Bn}
G(\xi,z)\Delta_h\varphi(z)\,d\nu(z)
=
\varphi(\xi).
\]
Hence
\[
\langle \Delta_h h,\varphi\rangle
=
\int_\Omega \varphi(\xi)\,d\nu(\xi)
=
\langle \chi_\Omega,\varphi\rangle.
\]
Therefore
\[
\Delta_h h=\chi_\Omega
\]
in the sense of distributions.

Inside \(\Omega\), the right-hand side equals the smooth function \(1\).
Elliptic regularity for \(\Delta_h\) implies that \(h\) is \(C^\infty\) in
\(\Omega\), and consequently
\[
\Delta_h h=1
\]
holds pointwise in \(\Omega\).
\end{proof}

\end{document}
